\chapter{Quantization of the Dirac Field}

\begin{remarkbox}
Scalar fields describe spin-zero quanta and are quantized with commutators. The
Dirac field describes spin-\(1/2\) quanta. Its consistent quantization requires
anticommutators, and its mode expansion naturally contains particle and
antiparticle operators with spin labels.
\end{remarkbox}

\section{Why Fermionic Quantization Is Different}

Chapters 2 and 3 quantized bosonic scalar fields. The ladder operators obeyed
commutation relations, and many-particle states were symmetric under exchange.
The Dirac field is different because its quanta are fermions. Fermionic states
must be antisymmetric under exchange, and no two identical fermions can occupy
the same one-particle state.

This difference is implemented algebraically by replacing commutators with
anticommutators. For operators \(\hat A\) and \(\hat B\), the anticommutator is
\[
    \{\hat A,\hat B\}
    = \hat A\hat B+\hat B\hat A.
\]
The quantized Dirac field will be built from creation and annihilation operators
whose anticommutation relations enforce Fermi statistics.

There is also a second important difference. The Dirac equation is first order
in time derivatives, unlike the Klein--Gordon equation. This changes the
canonical structure and makes the Dirac field look less like a collection of
ordinary oscillator coordinates. Nevertheless, the final Fock-space structure is
parallel to the scalar case: there are annihilation operators, creation
operators, a vacuum, particles, antiparticles, and number operators.

\section{Classical Dirac Field}

The Dirac field is a spinor field \(\psi(x)\). Its Dirac adjoint is
\[
    \bar\psi = \psi^\dagger\gamma^0.
\]
The free Dirac Lagrangian density is
\[
    \mathcal{L}
    = \bar\psi(i\gamma^\mu\partial_\mu-m)\psi.
\]
The equation of motion is the Dirac equation,
\[
    (i\gamma^\mu\partial_\mu-m)\psi(x)=0.
\]
The adjoint field satisfies
\[
    \bar\psi(x)(i\overleftarrow{\partial}_\mu\gamma^\mu+m)=0,
\]
where the arrow indicates that the derivative acts to the left.

The matrices \(\gamma^\mu\) obey the Clifford algebra
\[
    \{\gamma^\mu,\gamma^\nu\}=2\eta^{\mu\nu}.
\]
This algebra is what makes the Dirac equation compatible with the relativistic
mass-shell relation \(p^2=m^2\).

\section{Canonical Structure and Anticommutation Relations}

Because the Lagrangian is first order in \(\dot\psi\), the momentum conjugate to
\(\psi\) is proportional to \(\psi^\dagger\):
\[
    \pi_\psi
    = \frac{\partial\mathcal{L}}{\partial\dot\psi}
    = i\psi^\dagger.
\]
The quantum theory is defined by equal-time anticommutation relations,
\[
    \{\hat\psi_\alpha(t,\vect{x}),
      \hat\psi_\beta^\dagger(t,\vect{y})\}
    = \delta_{\alpha\beta}\delta^{(3)}(\vect{x}-\vect{y}),
\]
with
\[
    \{\hat\psi_\alpha(t,\vect{x}),
      \hat\psi_\beta(t,\vect{y})\}=0,
    \qquad
    \{\hat\psi_\alpha^\dagger(t,\vect{x}),
      \hat\psi_\beta^\dagger(t,\vect{y})\}=0.
\]
The spinor indices \(\alpha,\beta\) label the components of the Dirac spinor.

\begin{remarkbox}
The replacement of commutators by anticommutators is not a cosmetic change. It
is what produces fermionic Fock space, positive-energy particle interpretation,
and the Pauli exclusion principle.
\end{remarkbox}

\section{Mode Expansion of the Dirac Field}

The free Dirac equation has positive-energy spinor solutions \(u_s(\vect{p})\)
and negative-frequency spinor solutions \(v_s(\vect{p})\), where
\(s\) labels spin. A standard expansion is
\[
    \hat\psi(x)
    = \sum_s\int\frac{\dd^3p}{(2\pi)^3}
      \frac{1}{\sqrt{2E_{\vect{p}}}}
      \left(
          \hat b_s(\vect{p})u_s(\vect{p})e^{-iE_{\vect{p}}t+i\vect{p}\cdot\vect{x}}
          +
          \hat d_s^\dagger(\vect{p})v_s(\vect{p})e^{iE_{\vect{p}}t-i\vect{p}\cdot\vect{x}}
      \right),
\]
where
\[
    E_{\vect{p}}=\sqrt{\vect{p}^{\,2}+m^2}.
\]
The adjoint field is
\[
    \hat{\bar\psi}(x)
    = \sum_s\int\frac{\dd^3p}{(2\pi)^3}
      \frac{1}{\sqrt{2E_{\vect{p}}}}
      \left(
          \hat b_s^\dagger(\vect{p})\bar u_s(\vect{p})e^{iE_{\vect{p}}t-i\vect{p}\cdot\vect{x}}
          +
          \hat d_s(\vect{p})\bar v_s(\vect{p})e^{-iE_{\vect{p}}t+i\vect{p}\cdot\vect{x}}
      \right).
\]
The structure is analogous to the complex scalar field:
\[
    \hat\psi \sim \hat b + \hat d^\dagger,
    \qquad
    \hat{\bar\psi} \sim \hat b^\dagger + \hat d.
\]
The operators \(\hat b^\dagger\) create particles, while
\(\hat d^\dagger\) create antiparticles.

\section{Spinor Creation and Annihilation Operators}

The nonzero anticommutators are
\[
    \{\hat b_s(\vect{p}),\hat b_r^\dagger(\vect{q})\}
    = (2\pi)^3\delta_{sr}\delta^{(3)}(\vect{p}-\vect{q}),
\]
and
\[
    \{\hat d_s(\vect{p}),\hat d_r^\dagger(\vect{q})\}
    = (2\pi)^3\delta_{sr}\delta^{(3)}(\vect{p}-\vect{q}).
\]
All other anticommutators vanish. In particular,
\[
    \{\hat b_s(\vect{p}),\hat d_r(\vect{q})\}=0,
    \qquad
    \{\hat b_s(\vect{p}),\hat d_r^\dagger(\vect{q})\}=0.
\]
The label \(s\) is the spin label. For a massive Dirac particle there are two
spin states for the particle and two spin states for the antiparticle.

Because the creation operators anticommute, applying the same creation operator
twice gives zero:
\[
    \left(\hat b_s^\dagger(\vect{p})\right)^2=0,
    \qquad
    \left(\hat d_s^\dagger(\vect{p})\right)^2=0.
\]
This is the algebraic form of the Pauli exclusion principle.

\section{Particle and Antiparticle States}

The vacuum is defined by
\[
    \hat b_s(\vect{p})|0\rangle=0,
    \qquad
    \hat d_s(\vect{p})|0\rangle=0
    \qquad \text{for all }\vect{p},s.
\]
A one-particle state is
\[
    |\vect{p},s;b\rangle
    = \hat b_s^\dagger(\vect{p})|0\rangle.
\]
A one-antiparticle state is
\[
    |\vect{p},s;d\rangle
    = \hat d_s^\dagger(\vect{p})|0\rangle.
\]
For the electron field, the \(b^\dagger\)-states are electron states and the
\(d^\dagger\)-states are positron states. More generally, a Dirac field creates
and annihilates a charged fermion and its antiparticle.

A two-fermion state changes sign when the order of the creation operators is
exchanged:
\[
    \hat b_s^\dagger(\vect{p})\hat b_r^\dagger(\vect{q})|0\rangle
    = -\hat b_r^\dagger(\vect{q})\hat b_s^\dagger(\vect{p})|0\rangle.
\]
This antisymmetry is the defining feature of fermionic Fock space.

\section{Hamiltonian, Momentum, and Charge}

After normal ordering, the free Dirac Hamiltonian is
\[
    :\hat H:
    = \sum_s\int\frac{\dd^3p}{(2\pi)^3}\,
      E_{\vect{p}}
      \left(
          \hat b_s^\dagger(\vect{p})\hat b_s(\vect{p})
          +
          \hat d_s^\dagger(\vect{p})\hat d_s(\vect{p})
      \right).
\]
The momentum operator is
\[
    :\hat{\vect{P}}:
    = \sum_s\int\frac{\dd^3p}{(2\pi)^3}\,
      \vect{p}
      \left(
          \hat b_s^\dagger(\vect{p})\hat b_s(\vect{p})
          +
          \hat d_s^\dagger(\vect{p})\hat d_s(\vect{p})
      \right).
\]
Both particles and antiparticles carry positive energy.

For a Dirac field with unit charge convention, the normal-ordered charge has the
form
\[
    \hat Q
    = \sum_s\int\frac{\dd^3p}{(2\pi)^3}
      \left(
          \hat b_s^\dagger(\vect{p})\hat b_s(\vect{p})
          -
          \hat d_s^\dagger(\vect{p})\hat d_s(\vect{p})
      \right).
\]
For the electron field one often multiplies this operator by \(-e\), so that
particles created by \(\hat b^\dagger\) have charge \(-e\), and antiparticles
created by \(\hat d^\dagger\) have charge \(+e\).

\section{Why Anticommutators Give the Correct Energy Spectrum}

The negative-frequency part of the Dirac field is not interpreted as a tower of
negative-energy particles. After quantization, it is reinterpreted as the
creation of antiparticles with positive energy. The anticommutation relations are
essential for this reinterpretation to give a Hamiltonian bounded below.

Schematically, the Dirac Hamiltonian before normal ordering contains terms of
the form
\[
    \hat b^\dagger\hat b
    - \hat d\hat d^\dagger.
\]
Using anticommutators,
\[
    \hat d\hat d^\dagger
    = 1-\hat d^\dagger\hat d,
\]
so, up to an infinite constant that is removed by normal ordering, the
antiparticle contribution becomes positive:
\[
    \hat d^\dagger\hat d.
\]
This is the algebraic reason the free Dirac Hamiltonian counts both particles
and antiparticles with positive energy.

\section{Spinor Sums and Completeness}

The spinor wave functions are normalized so that they satisfy completeness
relations. A common convention is
\[
    \sum_s u_s(\vect{p})\bar u_s(\vect{p})
    = \gamma^\mu p_\mu + m,
\]
and
\[
    \sum_s v_s(\vect{p})\bar v_s(\vect{p})
    = \gamma^\mu p_\mu - m,
\]
where \(p^0=E_{\vect{p}}\). These identities are used constantly when computing
fermion propagators and scattering amplitudes.

The spinor sums play for the Dirac field a role analogous to polarization sums
for vector fields. They allow spin labels on internal lines to be summed in a
Lorentz-covariant way.

\section{The Dirac Propagator}

The Feynman propagator of the Dirac field is
\[
    S_F(x-y)
    = \langle0|T\hat\psi(x)\hat{\bar\psi}(y)|0\rangle.
\]
In momentum space it is
\[
    S_F(x-y)
    = \int\frac{\dd^4p}{(2\pi)^4}\,
      \frac{i(\gamma^\mu p_\mu+m)}{p^2-m^2+i\epsilon}
      e^{-ip\cdot(x-y)}.
\]
Equivalently, the propagator is the Green function of the Dirac operator:
\[
    (i\gamma^\mu\partial_\mu-m)S_F(x-y)
    = i\delta^{(4)}(x-y).
\]
Compared with the scalar propagator, the denominator is the same mass-shell
factor, but the numerator carries spinor structure.

\section{Chirality, Helicity, and the Massless Limit}

The Dirac field contains spin degrees of freedom. For massive fermions, the spin
label may be chosen in several equivalent ways. For massless fermions, helicity
and chirality become especially important.

The chirality matrix \(\gamma^5\) defines projection operators
\[
    P_L=\frac{1}{2}(1-\gamma^5),
    \qquad
    P_R=\frac{1}{2}(1+\gamma^5).
\]
They split a Dirac field into left- and right-chiral components:
\[
    \psi_L=P_L\psi,
    \qquad
    \psi_R=P_R\psi.
\]
For a massless Dirac field, the left- and right-chiral parts decouple. This will
be essential later in the electroweak theory, where left- and right-handed
fermions transform differently under gauge symmetries.

\section{Majorana Fields: Conceptual Overview}

A Dirac field describes a charged fermion and a distinct antiparticle. A
Majorana field is different: it is a fermionic field that is equal to its charge
conjugate in an appropriate sense. Informally, a Majorana fermion is its own
antiparticle.

The Majorana idea is possible for neutral fermions, not for fields carrying an
ordinary conserved electric charge. The full construction requires charge
conjugation matrices and reality conditions on spinors. For now, the important
contrast is simple:
\[
    \text{Dirac field: particle and antiparticle are distinct},
\]
while
\[
    \text{Majorana field: particle and antiparticle are identified}.
\]
Majorana fields will become relevant when discussing neutrino masses and
extensions of the Standard Model.

\begin{summarybox}
The Dirac field is the first fermionic quantum field in the course. Its
quantization uses anticommutators rather than commutators. The mode expansion
contains particle annihilation operators and antiparticle creation operators,
now with spin labels. Anticommutation gives fermionic Fock space, the Pauli
exclusion principle, and a Hamiltonian with positive particle and antiparticle
energies after normal ordering.
\end{summarybox}

\section{Chapter Checklist}

After reading this chapter, one should be able to:
\begin{itemize}
    \item write the free Dirac Lagrangian and Dirac equation;
    \item explain why fermionic fields are quantized with anticommutators;
    \item write the mode expansion of \(\hat\psi\) and \(\hat{\bar\psi}\);
    \item state the anticommutation relations of \(\hat b,\hat b^\dagger\) and
    \(\hat d,\hat d^\dagger\);
    \item construct fermion and antifermion states;
    \item explain how the Hamiltonian counts positive-energy particles and
    antiparticles;
    \item write the charge operator and interpret the sign convention;
    \item state the standard spinor completeness relations;
    \item write the Dirac propagator;
    \item distinguish Dirac and Majorana fermions at the conceptual level.
\end{itemize}
